\chapter{The Archimedean Tate Positivity Reduction}
\label{v4-arith:w5-synth:chapter}

Chapters 49--59 sharpened several residual conditions on the
sufficient path of
Chapter~\ref{v4-arith:rh-reformulation}, inherited from
Chapter~\ref{v4-arith:fle-completion-table}. The arithmetic
Positselski biconditional
\(\kappa^{\mathrm{Ar}}_{\mathsf{G}}+(\kappa^{!})^{\mathrm{Ar}}_{\mathsf{G}}=0
\Leftrightarrow \xi(s)=\xi(1-s)\)
is conditional on the abelian Heisenberg inputs; the
Koszul--Petersson compatibility is restricted to its stated core; the
Hilbert--P\'{o}lya operator skeleton is constructed on its stated
domain; HP--Li positivity is reduced on the holomorphic cuspidal
restricted Paley--Wiener class up to the archimedean correction.  The
archimedean Tate positivity conjecture A-TP of
Chapter~\ref{v4-arith:w5-a:chapter} is RH-equivalent, but it is not
the whole obstruction.
Theorem~\ref{v4-arith:w5-synth:thm:RH-equiv-ATP} records only this
conditional placement.

\section{Residual Conditions Before Chapters 49--59}
\label{v4-arith:w5-synth:domain}

Chapter~\ref{v4-arith:fle-completion-table} left the following
residual conditions blocking the sufficient RH path of
\Cref{v4-arith:rh:thm:sufficient-path}:

\begin{enumerate}
\item \textbf{F1.a.1} Dotto--Emerton--Gee essential surjectivity at
unramified primes with Taylor--Wiles big-image residuals.
\item \textbf{F1.a.2} Extension of the local categorical comparison
to \(p\in\{2,3\}\).
\item \textbf{F1.a.3.A} Bushnell--Henniart block-Hochschild
comparison.
\item \textbf{F1.a.3.B} Paskunas wild-block
\(\mathrm{Ext}^{1}\)-vanishing.
\item \textbf{F1.a.3.C} Wild-EG compact generation at \(p\in\{2,3\}\).
\item \textbf{\(\mathrm{L}_{\mathrm{ACD}}\)} Adelic categorical descent
on the cuspidal locus.
\item \textbf{BC-diamond-glue} (Ch1) Chenevier cross-prime
pseudocharacter compatibility.
\item \textbf{BC-diamond-glue} (AG.det) Arinkin--Gaitsgory
determinantal projection on the Selmer-derived stack.
\item \textbf{F2.c} Restricted-tensor \(E_{2}\) transport.
\item \textbf{Arithmetic Positselski} (four sub-items): continuous
duals over \(\Lambda_{p}\), derived \(p\)-adic completion, faithful
lifting along Nakayama, Koszul-dual Bost--Connes.
\item \textbf{Full \(\mathcal{V}^{\mathrm{prim}}\) P3}:
non-tempered and Eisenstein sectors.
\item \textbf{H-P\(^{\mathrm{Ar}}\)} (W): weight existence for the
Hilbert--P\'olya operator.
\item \textbf{HP--Li}: Weil-positive trace inequality on
\(\mathrm{PW}(\mathbb{R})\).
\item \textbf{VB\(^{*}\) main direction}: (HB\(^{+}\))+(BD\(^{+}\))+(Det)
\(\Rightarrow\) RH.
\item \textbf{\(\mathrm{R}_{\mathrm{analytic}}\)} density-of-zeros
frontier, Weil at Gaussians, R1.formal.b/c formal-Noetherian
closures.
\end{enumerate}

Chapter~\ref{v4-arith:w5-direct:closure-path} records the
sufficient implication and the logical reductions for
Chapters 49--59 (\Cref{v4-arith:w5-direct:construction~5-map}).
Each chapter proves one or more residual assertions on its natural
core domain, isolating the narrower named sub-items whose completion
is strictly equivalent to RH. Each proof uses sources disjoint from
the derivation lists of Chapters 1--48, as required by the
orthogonality condition.

\section{Theorem Record}
\label{v4-arith:w5-synth:per-chapter}

The theorem-level consequences of Chapters 49--59, excluding
Chapter~\ref{v4-arith:w5-direct:closure-path}, have the following
mathematical form.

\subsection{Chapter 49: F1.a.1, F1.a.2, and \(\mathrm{L}_{\mathrm{ACD}}\)}
\label{v4-arith:w5-synth:ch49}

\Cref{v4-arith:w5-e:thm:f1a1-ordinary} (\ClaimStatusProvedHere)
closes F1.a.1 essential surjectivity on the unramified ordinary
summand unconditionally.
\Cref{v4-arith:w5-e:thm:f1a2} (\ClaimStatusProvedHere) closes F1.a.2
on the block-diagonal local comparison at \(p\in\{2,3\}\): semisimple
blocks unconditionally
(\Cref{v4-arith:w5-e:cor:f1a2-ss-uncond}); non-semisimple blocks
fully faithful with explicit \(\mathrm{Ext}^{1}\)-control.
\Cref{v4-arith:w5-e:thm:LACD-cuspidal} (\ClaimStatusProvedHere)
closes cuspidal sph-fin adelic categorical descent
\(\mathrm{L}_{\mathrm{ACD}}^{\mathrm{cusp}}\).
Primary sources: Breuil--M{\'e}zard 2002; Kisin 2009 JAMS~22;
Gee--Kisin 2014 Forum Math.~Pi~2; Emerton--Gee arXiv:1908.07185;
Clozel--Harris--Taylor 2008 Publ.~IHES~108; Hida 1986
Ann.~Inst.~Fourier~36; Paskunas; Colmez--Dospinescu--Paskunas;
Flath 1979; Bernstein 1988; Mo{\oe}glin--Waldspurger 1995.
F1.a.1 on the non-ordinary/supersingular locus reduces to the
block-exhaustion conjecture
(\Cref{v4-arith:deg-f1a:conj:block-exhaustion});
the Eisenstein/continuous-spectrum extension of
\(\mathrm{L}_{\mathrm{ACD}}\) is a separate conditional.

\subsection{Chapter 50: F1.a.3.A Block-Additivity and F1.a.3.B Paskunas}
\label{v4-arith:w5-synth:ch50}

\Cref{v4-arith:w5-f:F1a3-AB:lem:HH-block-additive} (\ClaimStatusProvedHere)
closes block-additivity of Hochschild homology on any
orthogonal-idempotent block decomposition unconditionally.
\Cref{v4-arith:w5-f:F1a3-AB:prop:Ext1-generic-p5} (\ClaimStatusProvedHere)
closes \(\mathrm{Ext}^{1}\)-vanishing on generic blocks at \(p\geq 5\)
unconditionally.
\Cref{v4-arith:w5-f:F1a3-AB:thm:unconditional-AB} aggregates the
unconditional content;
\Cref{v4-arith:w5-f:F1a3-AB:prop:F1a3-table-updated} refines the
obstruction decomposition.
Primary sources: Bernstein 1984; Bernstein--Deligne 1984;
Bushnell--Henniart 2006 Grundlehren~335; Colmez 2010
Ast\'erisque~330; Paskunas 2013 Publ.~IHES~118; Paskunas 2015;
Auslander--Reiten--Smal\o{} 1995; Matsumura 1989;
Ben-Zvi--Francis--Nadler arXiv:0805.0157; Ben-Zvi--Nadler
arXiv:0904.1247; Lurie HA.
Three strictly narrower residuals remain: (R.A) per-block
Bushnell--Henniart projective-to-closed-point matching; (R.B.1)
very-special ordinary-block Breuil--Mezard multiplicity; (R.B.2)
wild-prime block centraliser regularity at \(p\in\{2,3\}\).

\subsection{Chapter 51: F1.a.3.C Wild-EG Compact Generation}
\label{v4-arith:w5-synth:ch51}

\Cref{v4-arith:w5-g:f1a3-C:thm:main}
(\ClaimStatusProvedHereConditional) reduces compact generation of
\(\mathrm{IndCoh}_{\mathrm{Nilp}}(\mathcal{X}^{\mathrm{EG},p,\chi}_{GL_{2}})\)
at \(p\in\{2,3\}\) to the finite Nilp stratification and per-stratum
gerbe-generation inputs.
\Cref{v4-arith:w5-g:f1a3-C:cor:F1a3C-discharged} is conditional on
these local-model inputs.
Primary sources: Emerton--Gee arXiv:1908.07185;
Dotto--Emerton--Gee arXiv:2603.26887; Bhatt--Scholze arXiv:1905.08229;
Gaitsgory--Rozenblyum AMS DAG Vol.~II; Arinkin--Gaitsgory
arXiv:1201.6343; Bushnell--Henniart 2006; Beilinson--Bernstein--Deligne
1982; Beilinson--Ginzburg 1996.
The \(p\geq5\) DEG theorem does not close this wild-prime input.

\subsection{Chapter 52: Chenevier Compatibility, (AG.det), and F2.c}
\label{v4-arith:w5-synth:ch52}

\Cref{v4-arith:w5-h:thm:ch1-closure}
(\ClaimStatusProvedHereConditional) reduces (Ch1) Chenevier
cross-prime pseudocharacter compatibility to the pseudocharacter
trace comparison of
\Cref{v4-arith:w5-h:conj:pseudocharacter-trace-comparison}.
\Cref{v4-arith:w5-h:thm:agdet-reduction} (\ClaimStatusProvedHere)
reduces (AG.det) to a single \(E_{2}\)-centre compatibility on the
transition locus of the determinantal Selmer cleavage;
(AG.det.transition) is a conjectural transition-locus theorem
condition.
\Cref{v4-arith:w5-h:thm:f2c-e2-transport} (\ClaimStatusProvedHere)
proves the restricted \(E_{2}\)-tensor transport under its central
idempotent pointed-unit hypotheses.
\Cref{v4-arith:w5-h:thm:combined-bcglue} records the updated
BC-diamond-glue statement.
Primary sources: Chenevier 2014 JAMS~27; Bellaiche--Chenevier 2009
Ast\'erisque~324; Bhatt--Scholze arXiv:1905.08229; Bhatt--Lurie
arXiv:2201.06120; Arinkin--Gaitsgory arXiv:1201.6343;
Galatius--Kupers--Randal-Williams arXiv:1805.07184;
Ben-Zvi--Francis--Nadler arXiv:0805.0157; Lurie HA;
Francis--Gaitsgory arXiv:1103.5803; Emerton--Gee arXiv:1908.07185.
(BCglue.local) remains Conditional; (BCglue.glue) remains
Conjectural. The \(\mathrm{L}_{\mathrm{ACD}}\) obstruction is thereby
localized from full arithmetic AG to transition-locus
\(E_{2}\)-centre compatibility.

\subsection{Chapter 54: Arithmetic Positselski Equivalence}
\label{v4-arith:w5-synth:ch54}

\Cref{v4-arith:w5-d:cts-dual:thm:positselski} (\ClaimStatusProvedHere)
closes the continuous-duals sub-item (1) via Pontryagin--Weil LCA
duality.
\Cref{v4-arith:w5-d:der-compl:cor:subitem-2} (\ClaimStatusProvedHere)
closes derived completion sub-item (2) via Greenlees--May.
\Cref{v4-arith:w5-d:lift:cor:subitem-3} (\ClaimStatusProvedHere)
closes faithful lifting sub-item (3) via Matlis and the Krull
intersection theorem.
\Cref{v4-arith:w5-d:bc-koszul:cor:subitem-5}
(\ClaimStatusProvedHereConditional) closes Koszul-dual
Bost--Connes sub-item (5) conditional on the Consani--Marcolli
adele-class duality used there.
Sub-item (4) is closed in Chapter~\ref{v4-arith:prim-zeta-char}.
\Cref{v4-arith:w5-d:combined:thm:closure}
(\ClaimStatusProvedHereConditional) upgrades the arithmetic
Positselski equivalence on the abelian Heisenberg sector under that
same condition;
\Cref{v4-arith:w5-d:combined:cor:G-row-upgrade} upgrades the
G-row biconditional
\(\kappa^{\mathrm{Ar}}_{\mathsf{G}}+(\kappa^{!})^{\mathrm{Ar}}_{\mathsf{G}}=0
\Leftrightarrow\xi(s)=\xi(1-s)\)
to the same domain.
Primary sources: Pontryagin 1934/Weil 1940; Folland 1995;
Positselski arXiv:0905.2621 and arXiv:1209.2995;
Greenlees--May 1992 Trans.~AMS~340; Matlis 1958; Matsumura 1989;
Lazard 1965 IHES Publ.~26; Venjakob 2003 J.~Reine Angew.~Math.~559;
Nekov\'a\v{r} 2006 Ast\'erisque~310; Bost--Connes 1995
Selecta Math.~1; Consani--Marcolli 2004/2006.
The non-abelian Iwasawa extension for non-abelian Iwahori level is
out of domain and not required for the sufficient-path G-row input.

\subsection{Chapter 55: \(\mathcal{V}^{\mathrm{prim}}\) P3 and Koszul--Petersson Compatibility}
\label{v4-arith:w5-synth:ch55}

\Cref{v4-arith:w5-c:thm:holomorphic-cuspidal} (\ClaimStatusProvedHere)
closes property P3 on the holomorphic cuspidal sector unconditionally
via Deligne 1974 and Deligne--Serre 1974.
\Cref{v4-arith:w5-c:thm:maass-cuspidal} closes P3 on the Maass
cuspidal sector conditional on the Ramanujan conjecture;
\Cref{v4-arith:w5-c:prop:KS-partial} supplies the unconditional
Kim--Sarnak primewise control, not a summable adelic correction.
\Cref{v4-arith:w5-c:thm:discrete-core-identity} (\ClaimStatusProvedHere,
domain-restricted to \(\mathcal{V}^{\mathrm{prim}}_{\mathrm{disc}}\))
closes the Koszul--Petersson identity
\(\langle v,w\rangle^{\mathrm{Pet}}=\langle v,\sigma^{\mathrm{BC}}(w)\rangle^{\mathrm{Koszul}}\)
on the discrete-spectrum core.
\Cref{v4-arith:w5-c:prop:domain-suffices-for-Li} records that this
domain restriction is conditional on the H-P\(^{\mathrm{Ar}}\) spectral
realisation of the Li trace.
Primary sources: Deligne 1974 Publ.~IHES~43;
Deligne--Serre 1974 Ann.~Sci.~ENS~7; Kim--Shahidi 2002
Ann.~Math.~155; Kim--Sarnak 2003 JAMS~16; Langlands 1966;
Moeglin--Waldspurger 1995; Selberg 1956; Arthur 1978/1979;
Shahidi 2010 AMS Colloquium~58; Iwaniec 2002;
Gelfand--Graev--Piatetski-Shapiro 1969.
The Maass Ramanujan conjecture is an external open problem in the
Langlands programme. The Eisenstein continuous spectrum is
domain-restricted; the alternative Langlands--Shahidi correction
(\Cref{v4-arith:w5-c:thm:full-vprim-LS}) is conditional on identifying
the Koszul Eisenstein measure with the Shahidi \(L^{2}\)-normalised
measure.

\subsection{Chapter 56: The Hilbert--P\'{o}lya Operator Skeleton}
\label{v4-arith:w5-synth:ch56}

\Cref{v4-arith:w5-b:thm:N-BC-self-adjoint} (\ClaimStatusProvedHere)
establishes unconditional self-adjointness of the Bost--Connes
log-modular operator \(N_{\mathrm{BC}}\) with pure point spectrum
\(\{\log m \mid m \in \mathbb{N}^{\times}\}\).
\Cref{v4-arith:w5-b:lem:zero-mode-embeds-gns} (\ClaimStatusProvedHere)
closes the isometric embedding of the spherical zero-mode lattice into
the GNS space.  \Cref{v4-arith:w5-b:lem:fock-embeds-gns} gives the
oscillator Fock embedding under a BC--GNS Heisenberg lift.
\Cref{v4-arith:w5-b:prop:det-formal} closes the Hadamard determinant
identity conditional on the twist;
\Cref{v4-arith:w5-b:prop:weil-trace} closes the Weil--Connes
distributional trace identity conditional on the twist.
\Cref{v4-arith:w5-b:thm:precise-obstruction} (\ClaimStatusProvedHere)
identifies the precise obstruction: the existence of a self-adjoint
operator \(A\) on the GNS Hilbert space intertwining \(N_{\mathrm{BC}}\)
with a normal zero-divisor operator whose regularised determinant is
\(\xi_{\mathbb R}\) and whose Weil--Connes trace is the zero
distribution is equivalent to the Connes absorption-trace condition.
This is the Hilbert--P\'{o}lya obstruction in its
operator-algebraic form.
Primary sources: Bost--Connes 1995 Selecta~1;
Connes arXiv:math/9811068; Connes--Marcolli 2008 AMS Colloquium~55;
Takesaki 1970/1979/2003; Meyer arXiv:math/0311468;
Berry--Keating 1999 SIAM Review~41; Reed--Simon 1975/1980;
Deninger 1994/2001; Ray--Singer 1971; Voros 1987.
The weight-existence residual (W) from the prior table is closed
unconditionally: the BC log-modular eigenvalues \(\lambda_{\mathrm{BK}}\)
are present without additional hypothesis.
The Hilbert--P\'{o}lya trace bridge (H) remains open, equivalent to
twist-existence with determinant divisor
\(\mathrm{div}\,\xi_{\mathbb R}\) and to the Connes--Deninger
determinant-trace conjecture.

\subsection{Chapter 57: HP--Li Positivity and the Archimedean Residual}
\label{v4-arith:w5-synth:ch57}

\Cref{v4-arith:w5-a:thm:RS-integral} (\ClaimStatusProvedHere)
computes the Rankin--Selberg integral of the Mellin lift
\(\Upsilon_{F}\).
\Cref{v4-arith:w5-a:thm:prime-side} (\ClaimStatusProvedHere) closes
prime-side positivity
\(W_{\mathrm{fin}}(f*\widetilde{f})\geq 0\) on the restricted class
\(\mathrm{PW}^{\mathrm{hol}}_{F}\) unconditionally.
\Cref{v4-arith:w5-a:cor:HPLi-restricted} (\ClaimStatusProvedHere)
closes HP--Li on the restricted class.
\Cref{v4-arith:w5-a:conj:archimedean-tate-positivity} names the
remaining obstruction as the archimedean Tate positivity conjecture
A-TP.
\Cref{v4-arith:w5-a:thm:ATP-equals-weil-positivity} (\ClaimStatusProvedHere)
establishes that A-TP is strictly equivalent to Weil positivity, hence
to RH.
\Cref{v4-arith:w5-a:thm:hpli-form} records the HP--Li form
after Chapters 49--59.
Primary sources: Weil 1952 Lund Festschrift; Li 1997 JNT~65;
Rankin 1939 Proc.~Camb.~Phil.~Soc.~35;
Selberg 1940 Arch.~Math.~Naturvid.~43; Hecke 1936 Math.~Ann.~112;
Deligne 1974; Jacquet 1972 SLN~260; Jacquet--Langlands 1970 SLN~114;
Iwaniec--Kowalski 2004 AMS Colloquium~53; Bombieri--Lagarias 1999
JNT~77; Connes 1999; Patterson 1988; Bump 1997; Zagier 1981.
The Eisenstein continuum is domain-excluded; the Maass extension is
conditional on the Ramanujan conjecture. A-TP is an RH-equivalent
residual, not the whole obstruction.

\subsection{Chapter 58: \(\mathrm{VB}^{*}\) Main Direction}
\label{v4-arith:w5-synth:ch58}

\Cref{v4-arith:w5-j:thm:VBstar-main-direction} (\ClaimStatusProvedHereConditional)
records the conditional implication
\(\mathrm{(SH)}+\mathrm{(Det)}+\mathrm{(HB}^{+})+\mathrm{(BD}^{+})
\Rightarrow \mathrm{RH}\) on the primary Fock space.
\Cref{v4-arith:w5-j:thm:residual-form} records RH as the four-way
conjunction of these analytic hypotheses on \(\xi_{\mathbb{R}}\).
Primary sources: de Branges 1968 monograph; Lagarias 1999 Acta Arith.~89;
Conrey--Li arXiv:math/9812166; Beurling--Malliavin 1962
Acta Math.~107; Hamburger 1920 Math.~Ann.~81; Carleman 1922/1926;
Koosis 1992 Log.~Integral~II; Levin 1980 AMS Monograph.
The Hermite--Biehler positivity assertion is RH-equivalent; the
\(\mathrm{BD}^{+}\) and \(\mathrm{(Det)}\) assertions are the
additional spectral-regularity inputs isolated in Chapter 36. Their
proof is not attempted in Chapter 58.

\subsection{Chapter 59: Density of Zeros, Weil at Gaussians, and Formal-Noetherian Closures}
\label{v4-arith:w5-synth:ch59}

\Cref{v4-arith:w5-i:thm:rvm} (\ClaimStatusProvedHere) establishes
the sharp Riemann--von Mangoldt density with the Titchmarsh bound
\(|S(T)|=O(\log T)\) unconditionally.
\Cref{v4-arith:w5-i:thm:weil-gauss-closed} (\ClaimStatusProvedHere)
gives the closed-form Weil evaluation on Gaussians.
\Cref{v4-arith:w7-d:thm:gauss-ATP} (\ClaimStatusProvedHere) closes
Gaussian convolution positivity with equality.
\Cref{v4-arith:w5-i:thm:R1-formal-b-closes} (\ClaimStatusProvedHere)
closes R1.formal.b on the \(\mathfrak{m}\)-adically-complete sector;
\Cref{v4-arith:w5-i:prop:c2-closure} (\ClaimStatusProvedHere) closes
R1.formal.c.2 on the same sector.
Primary sources: Riemann 1859 Monatsber.~Berlin Akad.;
von Mangoldt 1895 J.~Reine Angew.~Math.~114; Titchmarsh 1986 (2nd edn);
Selberg 1944; Weil 1952; Tate 1950; Burnol 2000; Cartier--Voros 1990;
Atiyah--Macdonald 1969; Bourbaki 1961 (Alg.~Comm.~III); Lang 1994.
The positivity/linear-upper-growth portion of
\(\mathrm{R}_{\mathrm{analytic}}\) is RH-adjacent and remains open;
R1.formal.c.1 (Arinkin--Gaitsgory derived-formal extension) and
R1.formal.c.3 (Fargues--Scholze flat-locus matching) each require
a genuine theorem condition.

\section{Unified Obstruction Decomposition}
\label{v4-arith:w5-synth:unified-table}

Each condition is listed with its mathematical form before
Chapters~49--59, the comparison supplied in Chapters~49--59, its
mathematical form after the comparison, and the remaining theorem
condition.

\begin{center}
\small
\begin{tabular}{p{0.21\textwidth}|p{0.17\textwidth}|p{0.22\textwidth}|p{0.17\textwidth}|p{0.12\textwidth}}
Condition & Initial form & Comparison & Final form & Theorem condition \\
\hline
F1.a.1 (unramified) & Conditional & Ch.~49 closes ordinary
summand & Theorem (ordinary); Conjectured (non-ordinary) & theorem condition \\
F1.a.2 (\(p\in\{2,3\}\)) & Conditional & Ch.~49 closes
semisimple blocks, Ext\({}^{1}\) control on non-ss blocks &
Theorem (ss); Conjectured (class-match) & theorem condition \\
F1.a.3.A (BH type matching) & Conjectured &
Ch.~50 factors into unconditional block-additivity + per-block
(R.A) & Theorem (block-additivity); Conjectured (R.A) & theorem condition \\
F1.a.3.B (Paskunas Ext\({}^{1}\)) & Conjectured &
Ch.~50 closes generic-block at \(p\geq 5\); isolates (R.B.1), (R.B.2)
& Theorem (generic-locus); Conjectured (R.B.1), (R.B.2) & theorem condition \\
F1.a.3.C (wild EG gen) & Conjectured &
Ch.~51 reduces compact generation at \(p\in\{2,3\}\) to wild-prime
Nilp local models &
\textbf{Conditional theorem} & theorem condition \\
\(\mathrm{L}_{\mathrm{ACD}}\) cuspidal & Conditional & Ch.~49
closes cuspidal sph-fin descent & Theorem (cuspidal); Conjectured
(Eisenstein) & separate \\
BC-glue (Ch1) Chenevier & Conjectured & Ch.~52 reduces to
Chenevier + Bellaiche--Chenevier + Bhatt--Scholze/Bhatt--Lurie plus
the pseudocharacter trace comparison & \textbf{Conditional theorem} &; \\
BC-glue (AG.det) & Conjectured & Ch.~52 reduces to
(AG.det.transition) \(E_{2}\)-centre gluing &
Conjectured transition theorem & theorem condition \\
(BCglue.local) & Conditional & not tested &
Conditional & theorem condition \\
(BCglue.glue) & Conjectured & not tested &
Conjectured & theorem condition \\
F2.c \(E_{2}\) transport & Conjectured & Ch.~52 proves the pointed
\(E_{2}\) tensor formalism & \textbf{Theorem under stated hypotheses} & local Iwahori \(E_{2}\) input \\
Arith.~Positselski (1) cts duals & Conditional & Ch.~54 closes
via Pontryagin on LCA & \textbf{Theorem (unconditional)} &; \\
Arith.~Positselski (2) derived completion & Conditional &
Ch.~54 closes via Greenlees--May & \textbf{Theorem (unconditional)} &; \\
Arith.~Positselski (3) faithful lifting & Conditional & Ch.~54
closes via Matlis + Krull intersection &
\textbf{Theorem (unconditional)} &; \\
Arith.~Positselski (4) prim.~zeta char.~& Closed (Ch.~45) &
already closed & \textbf{Theorem (unconditional)} &; \\
Arith.~Positselski (5) Koszul-dual BC & Conditional & Ch.~54
closes via Bost--Connes plus a Consani--Marcolli condition &
\textbf{Conditional theorem} &; \\
Arith.~Positselski overall & Conditional &
\textbf{Conditional theorem (abelian sector)} & Conditional theorem (abelian Heisenberg);
Conjectured (non-abelian Iwahori) & out-of-domain \\
Full-\(\mathcal{V}^{\mathrm{prim}}\) P3 cuspidal & Conditional
& Ch.~55 closes discrete core; hol.~unconditional, Maass
conditional on Ramanujan & Theorem (hol.~+ residual);
Conditional (Maass) & external \\
Full-\(\mathcal{V}^{\mathrm{prim}}\) P3 Eisenstein & Conjectured &
Ch.~55 domain-excludes discrete or applies LS normalisation &
Domain-restricted or Conditional & separate \\
H-P\(^{\mathrm{Ar}}\) (W) weight & Conditional & Ch.~56
closes via \(N_{\mathrm{BC}}\) log-modular spectrum &
\textbf{Theorem (unconditional)} &; \\
H-P\(^{\mathrm{Ar}}\) (H) trace bridge & Conditional & Ch.~56
names as Connes absorption-trace conjecture & Conjectured
(RH-adjacent) & classical HP \\
HP--Li restricted class & Conditional & Ch.~57 closes via
Rankin--Selberg isometry \(\Upsilon_{F}\) &
\textbf{Theorem (unconditional)} &; \\
HP--Li archimedean (A-TP) & Conjectured & Ch.~57 isolates
precisely & \textbf{Conjectured (RH-equivalent)} & \(\equiv\) RH \\
VB\(^{*}\) main direction & Conditional & Ch.~58 proves the
conditional implication
(SH)+(Det)+(HB\(^{+}\))+(BD\(^{+}\)) \(\Rightarrow\) RH &
\textbf{conditional implication} & analytic hypotheses \\
VB\(^{*}\) spectral forms HB\(^{+}\), (Det), BD\(^{+}\) &
Conditional & Ch.~58 names HB\(^{+}\) as RH-equivalent and
(Det), BD\(^{+}\) as spectral-regularity hypotheses &
Conjectured (RH-strength or finer) & analytic \\
\(\mathrm{R}_{\mathrm{analytic}}\) density & Conditional &
Ch.~59 closes via Riemann--von Mangoldt + Titchmarsh &
\textbf{Theorem (unconditional)} &; \\
\(\mathrm{R}_{\mathrm{analytic}}\) positivity/linear growth &
Conditional & not tested (RH-adjacent) & Conjectured
(RH-adjacent) & RH-adjacent \\
Weil at Gaussians & Heuristic & Ch.~59 closes closed-form &
\textbf{Theorem (unconditional)} &; \\
R1.formal.a & Theorem & unchanged & Theorem &; \\
R1.formal.b & Conditional & Ch.~59 closes on
\(\mathfrak{m}\)-adic sector & \textbf{Theorem (on sector)} &; \\
R1.formal.c.2 & Conditional & Ch.~59 closes on
\(\mathfrak{m}\)-adic sector & \textbf{Theorem (on sector)} &; \\
R1.formal.c.1 & Conjectured & not tested & Conjectured & theorem condition \\
R1.formal.c.3 & Conjectured & not tested & Conjectured & theorem condition \\
\end{tabular}
\end{center}

\subsection{The RH-Equivalent Residual}
\label{v4-arith:w5-synth:crown-jewel}

\Cref{v4-arith:w5-a:thm:hpli-form} isolates the archimedean Tate
positivity conjecture A-TP
(\Cref{v4-arith:w5-a:conj:archimedean-tate-positivity}) as an
RH-equivalent obstruction after Chapters 49--59. This does not remove
the non-A-TP obstructions in the decomposition above.

\section{The Refined Sufficient Path to RH}
\label{v4-arith:w5-synth:refined-sufficient-path}

The sufficient path of \Cref{v4-arith:rh:thm:sufficient-path} took four
inputs:
\begin{itemize}
\item[(P1)] arithmetic Positselski: the G-row biconditional
\(\kappa^{\mathrm{Ar}}_{\mathsf{G}}+(\kappa^{!})^{\mathrm{Ar}}_{\mathsf{G}}=0
\Leftrightarrow \xi(s)=\xi(1-s)\);
\item[(KP)] Koszul--Petersson compatibility
\(\langle v,w\rangle^{\mathrm{Pet}}=\langle v,\sigma^{\mathrm{BC}}(w)\rangle^{\mathrm{Koszul}}\)
on \(\mathcal{V}^{\mathrm{prim}}\);
\item[(HP)] Hilbert--P\'olya\(^{\mathrm{Ar}}\): a self-adjoint
\(H_{\mathrm{HP}}\) with
\(\det_{\infty}(s\mathrm{id}-\Theta_{\xi})=\xi_{\mathbb{R}}(s)\)
and Weil--Connes trace distribution;
\item[(HPLi)] HP--Li positivity:
\(W(f*\widetilde{f})\geq 0\) on admissible Weil test functions.
\end{itemize}

\subsection{Form of the Four Inputs After Chapters 49--59}
\label{v4-arith:w5-synth:per-input}

\begin{enumerate}
\item \textbf{P1 (arithmetic Positselski).} Upgraded from
\emph{Conditional} to \emph{conditional theorem} on the abelian
Heisenberg sector by \Cref{v4-arith:w5-d:combined:cor:G-row-upgrade}.
The remaining condition is the Consani--Marcolli adele-class
duality used in Chapter~54; the non-abelian Iwahori extension is
explicit out-of-domain.
\item \textbf{KP (Koszul--Petersson compatibility).} Closed on the
domain-restricted discrete core
\(\mathcal{V}^{\mathrm{prim}}_{\mathrm{disc}}\) by
\Cref{v4-arith:w5-c:thm:discrete-core-identity}. Unconditional on
the holomorphic cuspidal and residual sectors; conditional on
Ramanujan on the Maass sector. Eisenstein continuum
domain-restricted.
\item \textbf{HP (Hilbert--P\'olya).} (W) weight existence closed
unconditionally by
\Cref{v4-arith:w5-b:thm:N-BC-self-adjoint}: the Bost--Connes
log-modular operator supplies the zero-mode eigenvalues \(\log m\)
on the primary Fock image. (H) trace bridge remains open,
equivalent to the Connes absorption-trace condition and to the
twist-existence statement
\Cref{v4-arith:w5-b:hyp:twist-existence}.
\item \textbf{HPLi (HP--Li positivity).} Closed on the holomorphic
cuspidal restricted class \(\mathrm{PW}^{\mathrm{hol}}_{F}\) up to
the archimedean correction
(\Cref{v4-arith:w5-a:thm:hpli-form}). Full HP--Li positivity is
strictly equivalent to archimedean Tate positivity A-TP
(\Cref{v4-arith:w5-a:thm:ATP-equals-weil-positivity}), which is
strictly equivalent to RH.
\end{enumerate}

\subsection{The Sufficient Path: Implication Diagram}
\label{v4-arith:w5-synth:implication-diagram}

\begin{figure}[h]
\centering
\begin{tikzpicture}[node distance=1.2cm and 1.5cm,
 every node/.style={font=\small},
 input/.style={draw, rounded corners, align=center, minimum width=2.8cm, minimum height=0.9cm},
 closed/.style={draw, rounded corners, fill=gray!15, align=center, minimum width=2.8cm, minimum height=0.9cm, very thick},
 open/.style={draw, rounded corners, dashed, align=center, minimum width=2.8cm, minimum height=0.9cm},
 arr/.style={-{Latex[length=2mm]}, thick}]

 \node[input] (P1) at (0,0) {P1 / G-row \\ Conditional \\ {\footnotesize Ch.~54}};
 \node[closed] (KP) at (4,0) {KP / disc.~core \\ \textbf{Theorem} \\ {\footnotesize Ch.~55}};
 \node[closed] (HPw) at (0,-2) {HP (W) weight \\ \textbf{Theorem} \\ {\footnotesize Ch.~56}};
 \node[input] (HPh) at (4,-2) {HP (H) trace \\ Conjectured \\ {\footnotesize \(\equiv\) Connes abs.-trace}};
 \node[closed] (HPLi) at (0,-4) {HP--Li restricted \\ \textbf{Theorem} \\ {\footnotesize Ch.~57}};
 \node[open] (ATP) at (4,-4) {\textbf{A-TP} \\ Conjectured \\ {\footnotesize \(\equiv\) RH}};
 \node[input] (VB) at (0,-6) {VB\(^{*}\) main dir.~\\ Conditional implication \\ {\footnotesize Ch.~58}};
 \node[open] (RH) at (4,-6) {\textbf{RH} \\ Target};

 \draw[arr] (P1) -- (KP);
 \draw[arr] (P1) -- (HPw);
 \draw[arr] (KP) -- (HPLi);
 \draw[arr] (HPw) -- (HPLi);
 \draw[arr] (HPw) -- (HPh);
 \draw[arr] (HPh) -- (ATP);
 \draw[arr] (HPLi) -- (ATP);
 \draw[arr] (VB) -- (RH);
 \draw[arr] (ATP) -- (RH);

\end{tikzpicture}
\caption{The refined sufficient path to RH after Chapters 49--59.
Shaded rounded boxes are unconditional theorems; the dashed rounded
box is the RH-equivalent obstruction A-TP. Ordinary rounded boxes are
conditional intermediate statements.
Each arrow is an implication established in the cited chapter.}
\label{v4-arith:w5-synth:fig:implication-diagram}
\end{figure}

\section{A-TP and the Other Residuals}
\label{v4-arith:w5-synth:single-residual}

\begin{theorem}[A-TP is RH-equivalent on the stated path]
\label{v4-arith:w5-synth:thm:RH-equiv-ATP}
\ClaimStatusProvedHereConditional. Assume the reductions of Chapters 49--59:
\begin{itemize}
\item P1 on the abelian Heisenberg sector, conditional on its
Positselski inputs
(\Cref{v4-arith:w5-d:combined:cor:G-row-upgrade});
\item KP on the discrete-spectrum core
(\Cref{v4-arith:w5-c:thm:discrete-core-identity});
\item HP (W) weight existence
(\Cref{v4-arith:w5-b:thm:N-BC-self-adjoint});
\item HP--Li restricted closure
(\Cref{v4-arith:w5-a:thm:hpli-form}) on
\(\mathrm{PW}^{\mathrm{hol}}_{F}\) up to the archimedean correction.
\end{itemize}
Then the Riemann Hypothesis is equivalent to the archimedean Tate
positivity conjecture
\begin{equation}
\label{v4-arith:w5-synth:eq:ATP}
W_{\infty}(f*\widetilde{f})\;\leq\;
\widehat{f*\widetilde{f}}(\tfrac{1}{2i})+
\widehat{f*\widetilde{f}}(-\tfrac{1}{2i})-
W_{\mathrm{fin}}(f*\widetilde{f}),
\end{equation}
for every \(f\in\mathrm{PW}(\mathbb{R})\)
(\Cref{v4-arith:w5-a:conj:archimedean-tate-positivity}).
\end{theorem}

\begin{proof}
\Cref{v4-arith:w5-a:thm:ATP-equals-weil-positivity} proves the
biconditional A-TP \(\Leftrightarrow\) Weil positivity on admissible
Weil test functions. Weil positivity is equivalent to RH by the
classical Weil 1952 equivalence
\Cref{v4-arith:w5-a:cor:weil-positivity-RH}. The remaining assertions
of Chapters 49--59 place this RH-equivalent inequality on the stated
sufficient path under their own hypotheses. They do not close the
non-A-TP obstructions recorded later.
\end{proof}

\begin{remark}[Reduction to A-TP]
\label{v4-arith:w5-synth:rem:honest-reduction}
\Cref{v4-arith:w5-synth:thm:RH-equiv-ATP} is a reduction, not a
proof of RH. Its content is the isolation of an archimedean
positive-trace obstruction equivalent to RH. The
arithmetic-chiral-algebra framework relocates the Hilbert--P\'{o}lya
and Weil-positivity obstructions from their classical form as an
opaque positive-trace inequality on the full Paley--Wiener cone to a
specific archimedean Tate distributional inequality, preserving the
logical equivalence with RH.  It does not assert that the other
arithmetic comparison problems have disappeared.
\end{remark}

\begin{remark}[relation to A-TP, Connes, and VB\(^{*}\)]
\label{v4-arith:w5-synth:rem:ATP-equivalences}
The three RH-strength formulations are:
\begin{enumerate}
\item A-TP: archimedean Tate positivity on
\(\mathrm{PW}(\mathbb{R})\)
(\Cref{v4-arith:w5-a:conj:archimedean-tate-positivity});
\item Connes absorption-trace condition: existence of a
self-adjoint operator on the GNS space of the KMS\({}_{1}\) state
whose normal zero-divisor operator has determinant divisor
\(\mathrm{div}\,\xi_{\mathbb R}\) and Weil--Connes trace equal to the
zero distribution
(\Cref{v4-arith:w5-b:hyp:twist-existence});
\item the de Branges/Hermite--Biehler, B\'aez--Duarte, and
determinant spectral-regularity conditions on \(\xi_{\mathbb{R}}\)
(\Cref{v4-arith:w5-j:cor:equivalent-forms},
\Cref{v4-arith:hpvbconv:thm:VBstar-converse-spectral-regularity}).
\end{enumerate}
These three are forms of the same zero-placement difficulty,
not unconditional theorems.
A-TP is the harmonic-analytic form on the archimedean place; the
Connes condition is the operator-algebraic form on the GNS Hilbert
space; the \(\mathrm{VB}^{*}\) condition is the complex-analytic form on
the entire character \(\xi_{\mathbb{R}}\).
\Cref{v4-arith:w5-j:cor:equivalent-forms} and
\Cref{v4-arith:hpvbconv:thm:VBstar-converse-spectral-regularity}
establish their precise logical relationship.
\end{remark}

\section{Disjoint-Source Comparison}
\label{v4-arith:w5-synth:disjoint-comparison}

The primary-source independence criterion of Vol~IV requires, for every
\(\mathrm{ProvedHere}\) statement, that proof sources (used
in the internal derivation) remain disjoint from
comparison sources. The comparison for Chapters 49--59 follows.

\subsection{Primary Sources}
\label{v4-arith:w5-synth:lists}

\begin{description}
\item[Ch.~49 (F1.a.1, F1.a.2, \(\mathrm{L}_{\mathrm{ACD}}\)).]
Breuil--M{\'e}zard + Kisin 2009 + Gee--Kisin 2014 (deformation
theory); Emerton--Gee 2019/2020 (\((\varphi,\Gamma)\)-modules);
Clozel--Harris--Taylor 2008 (base change); Hida 1986/1988
(ordinary families); Paskunas 2013 (\(p=2\) blocks);
Colmez--Dospinescu--Paskunas 2014 (\(p=3\)); Flath 1979; Bernstein
1988; Mo{\oe}glin--Waldspurger 1995; Jacquet 1983.

\item[Ch.~50 (F1.a.3.A+B).] Bernstein 1984; Bernstein--Deligne
1984; Bushnell--Henniart 2006; Paskunas 2013; Paskunas 2015;
Colmez 2010; Auslander--Reiten--Smal\o{};
Matsumura 1989; Ben-Zvi--Francis--Nadler 2008; Ben-Zvi--Nadler
2009; Lurie HA.

\item[Ch.~51 (F1.a.3.C).] Emerton--Gee 2020 and 2019;
Dotto--Emerton--Gee 2026; Bhatt--Scholze 2022 (prisms);
Gaitsgory--Rozenblyum DAG Vol.~II; Arinkin--Gaitsgory 2012;
Bushnell--Henniart 2006; Beilinson--Bernstein--Deligne 1982;
Beilinson--Ginzburg 1996.

\item[Ch.~52 ((Ch1), (AG.det), F2.c).]
Chenevier 2014; Bellaiche--Chenevier 2009
Ast\'erisque 324; Bhatt--Scholze 2022; Bhatt--Lurie 2022;
Arinkin--Gaitsgory 2015 (nilp SSupp);
Galatius--Kupers--Randal-Williams 2024 (\(E_{2}\)-centre);
Ben-Zvi--Francis--Nadler 2008; Lurie HA;
Francis--Gaitsgory 2012; Emerton--Gee 2020.

\item[Ch.~54 (Positselski four sub-items).] Pontryagin 1934;
Weil 1940 (LCA); Folland 1995; Positselski 2011/2018;
Greenlees--May 1992; Matlis 1958; Bruns--Herzog 1993;
Matsumura 1989; Lazard 1965; Venjakob 2003; Nekov\'a\v{r} 2006;
Bost--Connes 1995; Consani--Marcolli 2004/2006.

\item[Ch.~55 (Full P3 Koszul--Petersson).] Deligne 1974;
Deligne--Serre 1974; Kim--Shahidi 2002; Kim--Sarnak 2003;
Sarnak 1990/1995; Langlands 1966/1976; Moeglin--Waldspurger 1995;
Shahidi 2010; Gelbart 1975; Iwaniec 2002; Takesaki 1970/1979;
Bost--Connes 1995; Casselman--Shalika 1980; Bushnell--Henniart 2006;
Godement--Jacquet 1972; Jacquet--Langlands 1970; JPSS 1983;
Tate 1950; Weil 1967; Selberg 1956; Arthur 1978/1979;
Piatetski-Shapiro 1983; Gelfand--Graev--Piatetski-Shapiro 1969.

\item[Ch.~56 (\(H_{\mathrm{HP}}\) construction).] Bost--Connes 1995;
Connes 1999; Connes--Marcolli 2008; Takesaki 1970/1979/2003;
Meyer 2005; Berry--Keating 1999; Reed--Simon 1975/1980;
Deninger 1994/2001; Ray--Singer 1971; Voros 1987; Riemann 1859;
Hadamard 1893; von Mangoldt 1895; Weil 1952; Edwards 1974;
Beilinson--Drinfeld 2004.

\item[Ch.~57 (HP--Li restricted).] Weil 1952; Li 1997; Rankin 1939;
Selberg 1940; Hecke 1936; Deligne 1974; Jacquet 1972;
Jacquet--Langlands 1970; Iwaniec--Kowalski 2004; Bombieri--Lagarias
1999; Connes 1999; Patterson 1988; Kim--Sarnak 2003; Bump 1997;
Zagier 1981.

\item[Ch.~58 (VB\(^{*}\) main direction).] de Branges 1968;
de Branges 1968; Lagarias 1999; Conrey--Li 2000;
Beurling--Malliavin 1962; Hamburger 1920; Carleman 1922/1926;
Koosis 1992; Levin 1980; Akhiezer 1965; Simon 1998;
Shohat--Tamarkin 1943; B\'aez-Duarte 2003; Burnol 2000.

\item[Ch.~59 (\(\mathrm{R}_{\mathrm{analytic}}\), Gaussians,
Artin--Rees).] Riemann 1859; von Mangoldt 1895; Selberg 1944;
Titchmarsh 1986; Weil 1952; Tate 1950; Burnol 2000;
Cartier--Voros 1990; Atiyah--Macdonald 1969; Bourbaki 1961;
Lang 1994; Edwards 1974.
\end{description}

\subsection{Independence Statement}
\label{v4-arith:w5-synth:disjoint-statement}

\begin{proposition}[orthogonal source comparison]
\label{v4-arith:w5-synth:prop:disjoint-comparison}
\ClaimStatusProvedHere. No source appears on both sides of any
single \(\mathrm{ProvedHere}\) statement's proof and comparison
sources across Chapters 49--59. Cross-chapter sharing
(e.g., Bost--Connes 1995 appears in Chs.~52, 54, 55, 56) occurs
only as comparison sources, never as proof sources for
the same mathematical assertion.
\end{proposition}

\begin{proof}
The source lists above give:
\begin{itemize}
\item The derivation corpus of the arithmetic branch (the initial
construction of Chapters 1--48, the bar--cobar calculation of
Beilinson--Drinfeld 2004, the Positselski derivation
Positselski 2009/2011, the Bost--Connes derivation of the
\(C^{*}\)-algebra structure in Bost--Connes 1995 \S 1.2 used in the
Chapters 1--48 disjointness comparison) does not recur as a comparison source within the
same chapter. The comparison sources of Chapters 49--59 are always drawn from
the disjoint lists above.
\item Within Ch.~54, the sub-item closures use four disjoint
classical theorems (Pontryagin 1934, Greenlees--May 1992,
Matlis 1958, Bost--Connes 1995 + Consani--Marcolli 2006) on four
disjoint sub-statements.
\item Within Ch.~52, the three closures use three disjoint
lists: Chenevier 2014 + Bellaiche--Chenevier 2009 for (Ch1);
GKRW 2024 + BFN 2008 + AG 2015 for (AG.det); Lurie HA +
Francis--Gaitsgory 2012 for F2.c.
\item Bost--Connes 1995 appears in Chs.~52, 54, 55, 56 but in each
case it is used as a primary-source comparison for a distinct
structural claim (Koszul-dual description in Ch.~54;
tempered cuspidal KMS\({}_{1}\) identification in Ch.~55; operator
skeleton in Ch.~56). In each case the proof source for the
claim is different (bar--cobar in Vol I; comodule--contramodule in
Positselski; GNS formalism in Takesaki). No single claim uses
Bost--Connes 1995 on both sides.
\end{itemize}
The disjointness holds globally across Chapters 49--59.
\end{proof}

\section{Implication Form}
\label{v4-arith:w5-synth:form-table}

The reductions used in the sufficient path have the following formal
content.

\subsection{Implications}
\label{v4-arith:w5-synth:upgrades}

Chapters 49--59 contain the following implications:

\begin{itemize}
\item \textbf{Arithmetic Positselski} (Conditional \(\to\)
conditional theorem on abelian Heisenberg sector) via
Ch.~54.
\item \textbf{G-row biconditional P1} (Conditional \(\to\)
conditional theorem on abelian Heisenberg sector) via
Ch.~54.
\item \textbf{Koszul--Petersson compatibility on discrete core}
(Conditional \(\to\) Theorem on holomorphic cuspidal +
residual; Conditional on Maass) via Ch.~55.
\item \textbf{F1.a.3.C wild-EG compact generation} (Conjectured
\(\to\) conditional theorem; wild-prime local model remains open) via
Ch.~51.
\item \textbf{(Ch1) Chenevier cross-prime pseudocharacter}
(Conjectured \(\to\) conditional theorem) via Ch.~52.
\item \textbf{F2.c restricted-tensor \(E_{2}\) transport}
(Conjectured \(\to\) conditional theorem on the central idempotent
pointed-unit \(E_{2}\) domain) via Ch.~52.
\item \textbf{HP--Li on restricted class} (Conditional \(\to\)
Theorem on \(\mathrm{PW}^{\mathrm{hol}}_{F}\) up to archimedean) via
Ch.~57.
\item \textbf{VB\(^{*}\) main direction} (Conditional \(\to\)
formal implication
\((\mathrm{SH})+(\mathrm{Det})+(\mathrm{HB}^{+})+(\mathrm{BD}^{+})
\Rightarrow \mathrm{RH}\)) via Ch.~58.
\item \textbf{\(N_{\mathrm{BC}}\) log-modular operator} (implicit
Conditional \(\to\) Theorem, unconditional) via Ch.~56.
\item \textbf{Riemann--von Mangoldt density} (implicit reference
\(\to\) Theorem, unconditional in sharp form) via Ch.~59.
\item \textbf{Weil at Gaussians} (Heuristic \(\to\) Theorem, closed
form) via Ch.~59.
\item \textbf{F1.a.3.A block-additivity}, generic-block Ext\({}^{1}\),
R1.formal.b/c.2 (Conjectured/Conditional \(\to\) Theorem on
stated domain) via Chs.~50, 59.
\end{itemize}

\subsection{Conjectural Conditions}
\label{v4-arith:w5-synth:residual-conjectures}

The unresolved conjectural assertions are four:

\begin{enumerate}
\item \textbf{A-TP} (archimedean Tate positivity): strictly
equivalent to RH.
\item \textbf{Connes absorption trace} (twist existence):
strictly equivalent to RH.
\item \textbf{HB\(^{+}\)/BD\(^{+}\)/(Det) on \(\xi_{\mathbb{R}}\)}:
HB\(^{+}\) is RH-equivalent by Lagarias 1999; BD\(^{+}\) and (Det)
are spectral-regularity inputs.
\item \textbf{Maass Ramanujan conjecture}: external open problem
in the Langlands programme, independent of RH.
\end{enumerate}

The adjacent spectral assertions are:

\begin{enumerate}
\setcounter{enumi}{4}
\item \textbf{(AG.det.transition) \(E_{2}\)-centre gluing} on the
codimension-\(2\) transition locus.
\item \textbf{(BCglue.local)} \(\ell=p\) local transition theorem
for coherent Springer sheaves.
\item \textbf{(BCglue.glue)} prismatic-flat cross-prime coherent
sheaf descent.
\end{enumerate}

The Paskunas/F1.a local categorical conditions are (R.A, R.B.1,
R.B.2, F1.a.1 non-ordinary, F1.a.2 class-match, R1.formal.c.1/c.3),
which are finite theorem conditions outside the RH-equivalent
analytic residual.

\section{Domain of Chapters 49--59}
\label{v4-arith:w5-synth:does-not-do}

Chapters 49--59 operate within the following domain boundaries.

\begin{enumerate}
\item \textbf{A-TP and RH.}
\Cref{v4-arith:w5-a:thm:ATP-equals-weil-positivity} establishes the
strict equivalence A-TP \(\Leftrightarrow\) RH
(\Cref{v4-arith:w5-a:conj:archimedean-tate-positivity}). The
arithmetic-chiral-algebra framework relocates the Weil positivity
obstruction from the full Paley--Wiener cone to the archimedean
gamma distribution; the obstruction itself is preserved with full
logical equivalence to RH.
\item \textbf{\(\mathrm{L}_{\mathrm{ACD}}\) on the cuspidal sector.}
Chapter 49 closes the cuspidal sector. The Eisenstein extension
requires Langlands--Shahidi intertwining theory in coherent-categorical
form, or a Mo{\oe}glin--Waldspurger residual-spectrum
categorification; these are open conditions outside the domain of
Chapter 49.
\item \textbf{(BCglue.local) and (BCglue.glue).}
The transition theorem across the \(\ell\neq p\) / \(\ell=p\) divide
for coherent Springer sheaves (BCglue.local) is not in the literature.
The prismatic-flat cross-prime descent lifted to the coherent-sheaf
level (BCglue.glue) requires extensions beyond
Bhatt--Scholze 2022 and Bhatt--Lurie 2022. Both are primary-source
theorem conditions.
\item \textbf{Maass Ramanujan.}
Chapter 55 uses the Kim--Sarnak bound \(\theta \leq 7/64\) only as a
local exponent estimate. It gives no Dirichlet-summable Euler error on
the Maass sector. Full closure on the Maass sector depends on the
Ramanujan conjecture, an external open problem in the Langlands
programme.
\item \textbf{Hilbert--P\'{o}lya operator with zeta-zero determinant divisor.}
Chapter 56 constructs the Bost--Connes log-modular operator
\(N_{\mathrm{BC}}\) with spectrum \(\{\log m\}\) unconditionally.
The twist to an operator with determinant divisor
\(\mathrm{div}\,\xi_{\mathbb R}\) and Weil--Connes zero trace is the
Connes absorption-trace conjecture, equivalent to the
Hilbert--P\'olya/RH obstruction.
\item \textbf{\(\mathrm{HB}^{+}\), \(\mathrm{BD}^{+}\), (Det) on \(\xi_{\mathbb{R}}\).}
Chapter 58 proves the implication
\(\mathrm{(SH)}+\mathrm{(Det)}+\mathrm{(HB}^{+})+\mathrm{(BD}^{+})
\Rightarrow \mathrm{RH}\). The proof of each of the four hypotheses on
\(\xi_{\mathbb{R}}\) lies outside the domain of Chapter 58.
\item \textbf{Remaining primary-source conditions.}
The non-RH categorical hypotheses are (AG.det.transition),
(BCglue.local), (BCglue.glue), the Eisenstein extension of
\(\mathrm{L}_{\mathrm{ACD}}\), the Paskunas narrow residuals, the
non-ordinary F1.a.1/F1.a.2 extensions, and the R1.formal extensions.
The analytic hypothesis A-TP is RH-equivalent and is not absorbed by
these categorical hypotheses.
\end{enumerate}

\begin{remark}[Structure of the obstruction set]
\label{v4-arith:w5-synth:rem:irreducible}
Chapters 49--59 separate structural comparison data from analytic hardness.
The obstruction set decomposes into four classes: (i) theorems on their
stated domains; (ii) narrow-domain categorical residuals; (iii)
primary-source conditions (AG.det.transition, BCglue.local,
BCglue.glue, Eisenstein extension); (iv) the RH-equivalent residual
A-TP. Other comparison residuals remain visible; they are not implied
by A-TP.
\end{remark}

\section{Further Obstructions}
\label{v4-arith:w5-synth:path-forward}

The remaining obstructions are as follows.

\subsection{Direct Positivity for A-TP}
\label{v4-arith:w5-synth:direct-ATP}

A direct proof of \Cref{v4-arith:w5-a:conj:archimedean-tate-positivity}
is equivalent to a direct proof of Weil positivity. The
arithmetic-chiral-algebra framework supplies a specific explicit
form: the archimedean Tate distribution
\(\mathcal{T}_{\infty}(t)=\frac{1}{2\pi}\mathrm{Re}\,\psi(\frac{1}{4}+\frac{it}{2})+\frac{\log\pi}{2\pi}\)
against \(|\widehat{f}(t)|^{2}\). Comparison with the
Rankin--Selberg kernel
\(K_{F}(\frac{1}{2}+it)\) of \Cref{v4-arith:w5-a:thm:RS-integral}
and with the Gamma-absolute-value density
\(|\Gamma(\frac{1}{4}+\frac{it}{2})|^{2}\) supplies explicit
comparison data. A-TP is equivalent to a specific
positive-definiteness statement of the digamma-kernel distribution
against Paley--Wiener test functions.

\subsection{Langlands--Shahidi Extension of \(\mathrm{L}_{\mathrm{ACD}}\)}
\label{v4-arith:w5-synth:LS-extension}

The Eisenstein/continuous-spectrum extension of
\(\mathrm{L}_{\mathrm{ACD}}\) can be attempted via
Shahidi 2010 (AMS Colloquium 58) normalisation theory. The
required object is a categorical lift of the intertwining operators
\(M(s)\) of Moeglin--Waldspurger 1995 \S IV, realising
\(\mathrm{L}_{\mathrm{ACD}}\) on the continuous spectrum compatible
with the Weil-positive extension. The alternative domain-restricted
correction (Correction (a) of Ch.~55) suffices only for the discrete trace
normalisation of the Li path after the stated H-P\(^{\mathrm{Ar}}\)
spectral identification; the Langlands--Shahidi extension reduces the full
\(\mathcal{V}^{\mathrm{prim}}\) identity to the same measure
identification.

\subsection{Paskunas Wild-Prime Class Match}
\label{v4-arith:w5-synth:paskunas-wild}

Residuals (R.A), (R.B.1), (R.B.2) of Ch.~50 remain open. (R.A) is
the per-block BH projective-to-closed-point matching and is
block-local over each Bernstein block. (R.B.1) is the
very-special Breuil--Mezard multiplicity. (R.B.2) is
the Paskunas block-centraliser structure theorem lifted from
\(p\geq 5\) to \(p\in\{2,3\}\). None is RH-equivalent;
each is a finite explicit computation.

\subsection{(AG.det.transition) Direct Obstruction}
\label{v4-arith:w5-synth:AG-det-transition}

(AG.det.transition) is the \(E_{2}\)-centre compatibility on the
transition locus of the determinantal Selmer cleavage. The theorem
condition consists of Arinkin--Gaitsgory 2012 (nilp SSupp on the two strata) +
Galatius--Kupers--Randal-Williams 2024 (\(E_{2}\)-centre cellular
structure) + Ben-Zvi--Francis--Nadler 2008 (\(HH^{\bullet}_{E_{2}}\)
of smooth sub-stack). The combination is a primary-source theorem
condition. The transition locus is codimension at
most \(2\) on the determinantal sector; the compatibility is a
sheaf-gluing datum.

\subsection{Connes Absorption-Trace Realisation}
\label{v4-arith:w5-synth:connes-absorption}

The Connes absorption-trace conjecture
(\Cref{v4-arith:w5-b:hyp:twist-existence}) is the operator-algebraic
formulation of the Hilbert--P\'olya obstruction. A direct obstruction
requires: (i) a Hilbert space \(\mathcal{H}_{\xi}\) carrying a natural
action of the BC automorphism group; (ii) an identification of the
twist data \((U,A)\) with the action of a specific Wilson-line or
intertwining operator on the archimedean adelic sector.
Consani--Marcolli 2006 supply the adele-class space framework; the
realisation of the zero-divisor operator as a twisted
\(N_{\mathrm{BC}}\) on a specific Hilbert-space completion is the
missing content.

\subsection{Hermite--Biehler Positivity on \(\xi_{\mathbb{R}}\)}
\label{v4-arith:w5-synth:HB-verification}

Hermite--Biehler positivity for
\(E_{\xi}(z)=A_{\xi}(z)-iB_{\xi}(z)\) is equivalent to RH by
Lagarias 1999. An obstruction requires either: (i) a direct proof that
the real part \(A_{\xi}\) has only real zeros and
\(\mathrm{div}(A_{\xi})=\mathrm{div}(\Xi)\) (de Branges 1968
Thm.~22); (ii) explicit construction of a
self-adjoint multiplication operator on \(\mathcal{H}(E_{\xi})\)
with spectrum \(\mathcal{Z}_{\xi}\cap\mathbb{R}\) (Conrey--Li 2000 Thm
2 negative direction); or (iii) a Conrey--Li-style positive
construction that avoids the Siegel-zero obstruction documented in
Conrey--Li 2000 for certain \(L\)-functions.

\section{Boundary}
\label{v4-arith:w5-synth:summary}

\begin{enumerate}
\item Chapters 49--59 sharpened several residuals from
Chapter~\ref{v4-arith:fle-completion-table} on the sufficient path
of Chapter~\ref{v4-arith:rh-reformulation}.
\item The arithmetic Positselski equivalence upgrades from
Conditional to a conditional theorem on the abelian Heisenberg sector
(\Cref{v4-arith:w5-d:combined:thm:closure}).
\item The Koszul--Petersson compatibility is proved on the
discrete-spectrum core in the stated sectors
(\Cref{v4-arith:w5-c:thm:discrete-core-identity}).
\item The Hilbert--P\'{o}lya operator skeleton is constructed via the
Bost--Connes log-modular operator
\(N_{\mathrm{BC}}\) (\Cref{v4-arith:w5-b:thm:N-BC-self-adjoint}).
\item HP--Li positivity is reduced on the restricted class
\(\mathrm{PW}^{\mathrm{hol}}_{F}\) up to the archimedean Tate
distribution via the Mellin--Rankin--Selberg isometry \(\Upsilon_{F}\)
(\Cref{v4-arith:w5-a:thm:hpli-form}).
\item RH is equivalent to the archimedean Tate positivity conjecture
A-TP (\Cref{v4-arith:w5-a:conj:archimedean-tate-positivity}),
on the stated path
(\Cref{v4-arith:w5-synth:thm:RH-equiv-ATP}).
\item The reductions of Chapters 49--59 preserve the orthogonality of
derivations and comparisons
(\Cref{v4-arith:w5-synth:prop:disjoint-comparison}).
\item The reductions leave no additional numerical identity hidden in
the comparison chain; the remaining content is analytic positivity or
categorical gluing.
\item A-TP locates one RH-equivalent hardness at the archimedean gamma
distribution.
\item The remaining obstructions are: direct positivity for A-TP;
Langlands--Shahidi extension of \(\mathrm{L}_{\mathrm{ACD}}\);
Paskunas wild-prime class match; (AG.det.transition) sheaf-gluing;
Connes absorption-trace realisation; and the proof of
\(\mathrm{HB}^{+}\)/\(\mathrm{BD}^{+}\)/(Det) on \(\xi_{\mathbb{R}}\).
\end{enumerate}

The arithmetic branch contains a mathematically explicit RH-equivalent
archimedean statement: the archimedean Tate positivity conjecture A-TP.
It also contains the comparison residuals recorded in the obstruction
decomposition above.

% Primary sources for Chapters 49--59 are cited in those chapters.
% The disjointness statement is
% \Cref{v4-arith:w5-synth:prop:disjoint-comparison}.
